\section{Motivation}

This thesis grew out of a practical project in which software was built from the ground up using AI coding agents, across two development phases that each embody a different philosophy of working with an agent.

The first phase followed a deliberately loose approach: an agent was asked to implement features from scratch with structural guidance that began minimally and was extended incrementally over the course of development as specific problems surfaced. The goal was to observe how an AI agent behaves when given substantial implementation freedom within a single-agent, conversational workflow --- and where exactly that freedom starts to cause problems, whether in the form of inconsistency, architectural drift, or small but consequential implementation mistakes.

What came out of that first phase shaped the second one directly. Rather than assuming that better software simply requires a more capable model, the second phase focused on something different: the architectural rules, review steps, and development discipline imposed on the agents themselves. A more structured methodology was introduced for this purpose, built around role-specialized agents working together under a defined process rather than one agent working through a single, linear thread of tasks. The specific tools, agent roles, and artifacts used to realize this structure are introduced in Chapter 3.

The motivation behind this thesis is therefore twofold. First, it aims to document and compare, in concrete terms, what changes when moving from an unstructured, single-agent approach to a structured, multi-agent approach in AI-assisted software development. Second, it aims to distill from that comparison a set of best practices for human-agent collaboration that hold up beyond the specific tools used here, and that can inform how developers approach this kind of work more generally.

